% ===========================================================================
%  Bitfeld einer Gleitkomma-Darstellung
%  Thema 13: Zahlendarstellungen und Kodierungen
%
%  Aufruf:  \Bitfeld[<Schlüssel>=<Wert>, ...]
%
%    vz        Anzahl der Vorzeichenbits               (Standard: 1)
%    exp       Anzahl der Exponentbits                 (Standard: 8)
%    mant      Anzahl der Mantissenbits                (Standard: 23)
%    width     Gesamtbreite des Balkens in cm          (Standard: aus cell)
%    cell      Breite einer Bit-Zelle in cm            (Standard: 0.4)
%    height    Höhe des Balkens in cm                  (Standard: 0.72)
%    explabel  Feldname des Exponenten                 (Standard: Exponent)
%    mantlabel Feldname der Mantisse                   (Standard: Mantisse)
%
%  Im Balken steht das Massband der Bitpositionen (15 … 0), darunter benennen
%  Klammern die drei Felder samt ihrer Breite.
%
%  Das Standardformat ist IEEE 754 einfach (1/8/23), also 32 Bit.
%  Für das 16-Bit-Format aus Übung 1.32:  \Bitfeld[exp=5, mant=10, width=13.2]
% ===========================================================================
\pgfkeys{
  /bitfeld/.is family,
  /bitfeld/.cd,
  vz/.store in        = \bfvz,
  exp/.store in       = \bfexp,
  mant/.store in      = \bfmant,
  width/.store in     = \bfwidth,
  cell/.store in      = \bfcell,
  height/.store in    = \bfheight,
  explabel/.store in  = \bfexplabel,
  mantlabel/.store in = \bfmantlabel,
  % Standardwerte
  vz        = 1,
  exp       = 8,
  mant      = 23,
  width     = 0,
  cell      = 0.4,
  height    = 0.72,
  explabel  = {Exponent},
  mantlabel = {Mantisse},
}

\newcommand{\Bitfeld}[1][]{%
  \begingroup
  \pgfkeys{/bitfeld/.cd, #1}%
  % Anzahl Bits sowie die Feldgrenzen in Bit-Einheiten
  \pgfmathtruncatemacro{\bfn}{\bfvz+\bfexp+\bfmant}%
  \pgfmathtruncatemacro{\bfnlast}{\bfn-1}%
  \pgfmathtruncatemacro{\bfa}{\bfvz}%          Grenze Vorzeichen/Exponent
  \pgfmathtruncatemacro{\bfb}{\bfvz+\bfexp}%   Grenze Exponent/Mantisse
  % Zellbreite: entweder vorgegeben oder aus der Gesamtbreite berechnet
  \pgfmathsetmacro{\bfcell}{(\bfwidth>0.1) ? \bfwidth/\bfn : \bfcell}%
  % dieselben Grenzen als x-Koordinaten
  \pgfmathsetmacro{\bfxA}{\bfa*\bfcell}%
  \pgfmathsetmacro{\bfxB}{\bfb*\bfcell}%
  \pgfmathsetmacro{\bfxC}{\bfn*\bfcell}%
  \begin{tikzpicture}[x=1cm, y=1cm]
    % --- Felder hinterlegen ---------------------------------------------
    \fill[Crimson!10]     (0,0) rectangle (\bfxA,\bfheight);
    \fill[RoyalBlue!12]   (\bfxA,0) rectangle (\bfxB,\bfheight);
    \fill[ForestGreen!12] (\bfxB,0) rectangle (\bfxC,\bfheight);
    % --- Massband der Bitpositionen ------------------------------------
    \foreach \i in {0,...,\bfnlast}{
      \node[font=\scriptsize\ttfamily, text=black!45, inner sep=0pt]
        at ({(\i+0.5)*\bfcell},{0.5*\bfheight}) {\pgfmathparse{int(\bfnlast-\i)}\pgfmathresult};
    }
    % --- feine Zellgrenzen ---------------------------------------------
    \foreach \i in {1,...,\bfn}{
      \draw[black!18, line width=0.4pt]
        ({\i*\bfcell},0) -- ({\i*\bfcell},\bfheight);
    }
    % --- Rahmen und Feldgrenzen ----------------------------------------
    \draw[black!75, line width=0.9pt] (0,0) rectangle (\bfxC,\bfheight);
    \draw[black!75, line width=0.9pt] (\bfxA,0) -- (\bfxA,\bfheight);
    \draw[black!75, line width=0.9pt] (\bfxB,0) -- (\bfxB,\bfheight);
    % --- Feldnamen mit Klammer -----------------------------------------
    \def\bflabely{-0.46}%
    \draw[decorate, decoration={brace, amplitude=4pt, mirror}, black!55, line width=0.5pt]
      ([yshift=-1pt]0,0) -- ([yshift=-1pt]\bfxA,0);
    \node[font=\footnotesize, text=Crimson!80!black, align=center, inner sep=0pt]
      at ({0.5*\bfxA},\bflabely) {Vorzeichen\\[-4pt]{\scriptsize\bfvz\ Bit}};
    \draw[decorate, decoration={brace, amplitude=4pt, mirror}, black!55, line width=0.5pt]
      ([yshift=-1pt]\bfxA,0) -- ([yshift=-1pt]\bfxB,0);
    \node[font=\footnotesize, text=RoyalBlue!80!black, align=center, inner sep=0pt]
      at ({0.5*(\bfxA+\bfxB)},\bflabely) {\bfexplabel\\[-4pt]{\scriptsize\bfexp\ Bit}};
    \draw[decorate, decoration={brace, amplitude=4pt, mirror}, black!55, line width=0.5pt]
      ([yshift=-1pt]\bfxB,0) -- ([yshift=-1pt]\bfxC,0);
    \node[font=\footnotesize, text=ForestGreen!70!black, align=center, inner sep=0pt]
      at ({0.5*(\bfxB+\bfxC)},\bflabely) {\bfmantlabel\\[-4pt]{\scriptsize\bfmant\ Bit}};
  \end{tikzpicture}%
  \endgroup
}
